%! Simplifying Rational Expressions — Unit 4, Lesson 4.1 — Guided Notes (Answer Key)
\documentclass[10pt]{article}
\usepackage{algebra2-article}
\usepackage{algebra2-key}   % pulls in -boxes; do NOT also load -boxes
\newcommand{\TallMath}[1]{$\displaystyle #1\rule[-1.4em]{0pt}{3.2em}$}
\newcommand{\lgrid}[4]{%
  \draw[step=1, linegray!40, very thin] (#1,#3) grid (#2,#4);
  \draw[->, semithick, charcoal] (#1,0)--(#2,0) node[below right, font=\scriptsize]{$x$};
  \draw[->, semithick, charcoal] (0,#3)--(0,#4) node[left, font=\scriptsize]{$y$};}
% Vocab answer for the key: bold forest term, then the filled definition on a write-line.
% The leading and trailing \par are required: \ansline ends with \dotfill but does NOT end the
% paragraph, so without them each term label gets pulled onto the previous answer's line.
\newcommand{\vocabans}[2]{%
  \par\noindent\textbf{\textcolor{forest}{#1:}}\\[1pt]\ansline{#2}\par}

\begin{document}

\pageheader{Unit 4, Lesson 4.1}{Guided Notes --- Answer Key}
\namedateperiod

\vspace{0.08in}

\begin{objectivebox}
\textbf{By the end of this lesson, I will be able to\dots}
\begin{itemize}\setlength{\itemsep}{2pt}
  \item factor a rational expression completely and write it in \emph{simplest form} by dividing out common \emph{factors} --- never terms;
  \item list every \emph{domain restriction} from the \emph{original} denominator, including the ones that cancel;
  \item \emph{justify} that two rational expressions are equivalent, and say exactly which inputs they disagree on.
\end{itemize}
\end{objectivebox}

\vspace{0.05in}

\begin{vocabbox}
\small
Fill in each term as we build it together.
\par\vspace{2pt}   % \par is required: \vocabans's \noindent is a no-op mid-paragraph
\vocabans{Rational expression}{a quotient of two polynomials, $\frac{P(x)}{Q(x)}$, where $Q(x)\neq0$}
\vocabans{Simplest form (lowest terms)}{the form in which the numerator and denominator share no common factor except $1$ --- with the original restrictions still attached}
\vocabans{Factor vs.\ term}{factors are \emph{multiplied}; terms are \emph{added or subtracted}. Only common \emph{factors} may be divided out}
\vocabans{Equivalent rational expressions}{expressions that give the same value at every input for which \emph{both} are defined}
\vocabans{Opposite binomials}{two binomials with every sign flipped, such as $x-5$ and $5-x$; one is $-1$ times the other, so their quotient is $-1$}
\end{vocabbox}

\begin{hookbox}
Two expressions, one table. Let $A(x)=\dfrac{x^2-x-6}{x^2-9}$ and $B(x)=\dfrac{x+2}{x+3}$.
\begin{center}
\renewcommand{\arraystretch}{1.35}
\begin{tabular}{|l|c|c|c|c|}
\hline
$x$ & $0$ & $1$ & $4$ & $3$ \\ \hline
$A(x)$ & \ans{$\frac23$} & \ans{$\frac34$} & \ans{$\frac67$} & \ans{undef.} \\ \hline
$B(x)$ & \ans{$\frac23$} & \ans{$\frac34$} & \ans{$\frac67$} & \ans{$\frac56$} \\ \hline
\end{tabular}
\end{center}
\begin{itemize}\setlength{\itemsep}{1pt}
  \item Where do $A$ and $B$ agree? \ansline{everywhere they are both defined --- at $x=0,1,4$ and at every other allowed input}
  \item There is exactly one column where they do \emph{not} agree. Which one, and what goes wrong there? \ansline{at $x=3$: $A(3)=\frac{0}{0}$ is undefined (the denominator $x^2-9$ is zero), while $B(3)=\frac56$}
\end{itemize}
$B$ is what you get when you factor $A$ and cancel. It is \emph{almost} the same expression --- and the
whole lesson lives in that word \emph{almost}: \textbf{cancelling changes how an expression looks, but
it never gives back an input that was already thrown away.}
\end{hookbox}

\begin{notesbox}{1. Rational Expressions and Their Restrictions}
A \textbf{rational expression} is a \emph{quotient of two polynomials} --- the same object as a
rational function, just written without the ``$f(x)=$.'' Because division by zero is undefined, every
input making the denominator zero is a \textbf{domain restriction} (an \textbf{excluded value}).

\vspace{3pt}
\textbf{The procedure from Lesson 4.0, unchanged:} set the \ans{denominator} equal to \ans{0},
solve, and exclude every answer.
\begin{itemize}[leftmargin=1.5em, itemsep=3pt]
  \item For the anchor $A(x)=\dfrac{x^2-x-6}{x^2-9}$: solve $x^2-9=0$. \; Factored: \ans{$(x-3)(x+3)$},
        \; so the excluded values are $x=$ \ans{$3,\ -3$}.
  \item \textbf{Careful:} the \emph{numerator} is allowed to be zero --- that gives a \textbf{zero} of
        the expression, not a restriction.
\end{itemize}

\vspace{3pt}
\textbf{Find the restrictions \emph{before} you simplify.} That is not advice about tidiness --- it is
the only way to get them all, as the next two pages show.
\end{notesbox}

\begin{notesbox}{2. Why Cancelling Works --- and What It Is Not}
Cancelling is not crossing out matching symbols. It is one rule:
\[
  \frac{a\cdot c}{b\cdot c} \;=\; \frac{a}{b}\cdot\frac{c}{c} \;=\; \frac{a}{b}\cdot 1
  \;=\; \frac{a}{b}, \qquad b\neq0,\ c\neq0 .
\]
You are dividing out a factor equal to \ans{1}. This works \emph{only} when $c$ is a
\textbf{factor} --- something multiplied by everything else --- and only when $c\neq0$, which is
exactly why the restrictions matter.

\vspace{4pt}
\textbf{Factors are multiplied. Terms are added or subtracted.}
\begin{center}
\renewcommand{\arraystretch}{1.3}
\begin{tabularx}{\linewidth}{@{}l X X@{}}
\hline
 & \textbf{You may divide out} & \textbf{You may never divide out} \\
\hline
Name & a common \ans{factor} & a common \ans{term} \\
Example & $\dfrac{6(x+1)}{6(x-4)}=\dfrac{x+1}{x-4}$ & $\dfrac{6+x}{6-4}$ \; --- nothing cancels \\
\hline
\end{tabularx}
\end{center}

\vspace{2pt}
\textbf{Kill the error with a number.} A classmate claims $\dfrac{x+3}{3}=x$. Test it at $x=6$: the
left side $\dfrac{6+3}{3}$ equals \ans{3}, and the right side equals \ans{6}. They
disagree, so the ``rule'' is \ans{false}. In $\dfrac{x+3}{3}$ the $3$ on top is a \ans{term}
of the numerator, not a factor of it.
\end{notesbox}

\begin{notesbox}{3. The Four Steps}
\begin{enumerate}[leftmargin=1.9em, itemsep=2pt]
  \item \textbf{Factor} the numerator and the denominator \emph{completely} (GCF first --- the Unit 3 toolkit).
  \item \textbf{List the restrictions} from the \textbf{original} denominator's factors. \emph{Every} one.
  \item \textbf{Divide out} the factors the top and bottom share.
  \item \textbf{Write the simplest form}, and carry the restrictions along with it.
\end{enumerate}

\vspace{4pt}
\textbf{The anchor, worked.} $A(x)=\dfrac{x^2-x-6}{x^2-9}$
\begin{itemize}[leftmargin=1.5em, itemsep=3pt]
  \item Numerator factors as \ans{$(x-3)(x+2)$}; \; denominator factors as \ans{$(x-3)(x+3)$}.
  \item Restrictions, from the \emph{original} bottom: $x\neq$ \ans{$3,\ -3$}.
  \item The shared factor is \ans{$(x-3)$}. Divide it out: simplest form \ans{$\frac{x+2}{x+3}$}.
\end{itemize}

\vspace{3pt}
Now notice what just happened: the factor $(x-3)$ cancelled, but $x=3$ \textbf{is still excluded}. The
simplified expression looks perfectly happy at $x=3$; the \emph{original} is not, and the original is
what the question was about. In Lesson~4.0 you read that single missing point off a graph and called
it a \ans{hole}.

\vspace{5pt}
\textbf{Monomial factors work the same way.} Simplify $\dfrac{12x^3y}{18x^5y^2}$.
\begin{itemize}[leftmargin=1.5em, itemsep=3pt]
  \item Restrictions: the denominator is zero when $x=0$ or $y=0$, so $x\neq$ \ans{0} and
        $y\neq$ \ans{0}.
  \item Divide out the common numerical factor \ans{6} and the shared powers of $x$ and $y$:
        the answer is \ans{$\frac{2}{3x^2y}$}.
\end{itemize}

\vspace{4pt}
\textbf{One more, with a GCF.} Simplify $\dfrac{3x^2-12}{x^2+x-6}$.
\begin{itemize}[leftmargin=1.5em, itemsep=3pt]
  \item Factored numerator: \ans{$3(x-2)(x+2)$} \quad Factored denominator: \ans{$(x+3)(x-2)$}
  \item Restrictions: $x\neq$ \ans{$2,\ -3$}. \quad Simplest form: \ans{$\frac{3(x+2)}{x+3}$}.
\end{itemize}
\end{notesbox}

\begin{notesbox}{4. Opposite Binomials: the Hidden Factor of $-1$}
$x-5$ and $5-x$ are not the same, so they do not cancel --- but they are \textbf{opposites}, and
opposites divide to give $-1$:
\[
  \frac{5-x}{x-5} \;=\; \frac{-1(x-5)}{x-5} .
\]
So $\dfrac{5-x}{x-5}$ simplifies to \ans{$-1$}, and in general $\dfrac{a-b}{b-a}$ equals
\ans{$-1$}.
\begin{itemize}[leftmargin=1.5em, itemsep=3pt]
  \item \textbf{How to spot it:} two binomials with the same terms and \emph{every} sign flipped.
        Factor $-1$ out of one of them and they match.
  \item \textbf{Not the same situation:} $x+3$ and $3+x$ \emph{are} equal (addition can be reordered),
        so they cancel to \ans{1}. Only \emph{subtraction} order matters.
  \item \textbf{Worked:} $\dfrac{x^2-16}{4-x}=\dfrac{(x-4)(x+4)}{-(x-4)}$, which simplifies to
        \ans{$-(x+4)$}, \; for $x\neq$ \ans{4}.
\end{itemize}
\end{notesbox}

\begin{notesbox}{5. Equivalence --- and the Graph It Describes}
Two rational expressions are \textbf{equivalent} when they produce the same value at every input
\emph{both} are defined for. Simplifying always produces an equivalent expression --- but the two forms
can have different \emph{domains}, and \textbf{the restriction list belongs to the original}.

\vspace{3pt}
Here is the graph of the anchor $A(x)=\dfrac{x^2-x-6}{x^2-9}$, whose simplified form is
$\dfrac{x+2}{x+3}$ with $x\neq3$.
\begin{center}
\begin{tikzpicture}[scale=0.48]
  \lgrid{-8}{6}{-4}{5}
  \draw[navy, dashed, semithick] (-3,-4)--(-3,5);
  \draw[navy, dashed, semithick] (-8,1)--(6,1);
  \node[navy, font=\scriptsize, above left] at (-3,5) {$x=-3$};
  \node[navy, font=\scriptsize, below] at (-6.4,1) {$y=1$};
  \begin{scope}
    \clip (-8,-4) rectangle (6,5);
    \draw[very thick, forest, domain=-8:-3.25, samples=100, smooth]
      plot (\x,{((\x)+2)/((\x)+3)});
    \draw[very thick, forest, domain=-2.8:6, samples=100, smooth]
      plot (\x,{((\x)+2)/((\x)+3)});
  \end{scope}
  \draw[forest, very thick, fill=white] (3,0.833) circle (4.2pt);
  \fill[charcoal] (-2,0) circle (2.6pt);
  \node[charcoal, font=\scriptsize, below] at (-1.2,-0.35) {$(-2,0)$};
  \node[charcoal, font=\scriptsize] at (4.0,3.4) {the hole};
  \draw[charcoal, thin, ->] (3.7,3.0) -- (3.15,1.1);
\end{tikzpicture}
\end{center}
\begin{enumerate}[leftmargin=1.7em, itemsep=3pt]
  \item The two excluded values do \emph{different} jobs. At $x=-3$ the factor \emph{stays}, giving a
        \ans{vertical asymptote}; at $x=3$ the factor \emph{cancels}, giving a \ans{hole}.
  \item The hole's height comes from the \emph{simplified} form: substituting $x=3$ into
        $\dfrac{x+2}{x+3}$ gives \ans{$\frac56$}, so the hole sits at \ans{$\left(3,\frac56\right)$}.
  \item \textbf{Which form should you write?} The simplified one --- \emph{with} the restrictions
        attached. The simplified form tells you the shape of the graph; the restrictions tell you which
        points were punched out of it.
\end{enumerate}

\vspace{3pt}
\textbf{Checking equivalence two ways.} Are $\dfrac{x^2-4x}{x^2-16}$ and $\dfrac{x}{x+4}$ equivalent?
\begin{itemize}[leftmargin=1.5em, itemsep=3pt]
  \item \emph{By factoring:} $\dfrac{x(x-4)}{(x-4)(x+4)}$ simplifies to \ans{$\frac{x}{x+4}$} --- yes.
  \item \emph{By testing an input:} at $x=1$, the first gives $\dfrac{-3}{-15}=$ \ans{$\frac15$} and the
        second gives \ans{$\frac15$}. One matching test cannot \emph{prove} equivalence, but one failed
        test \emph{disproves} it --- which is exactly how the $x=6$ test settled the argument in
        Section~2.
  \item \emph{The catch:} the two still disagree at $x=$ \ans{4}, where the first expression is
        \ans{undefined ($\frac00$)} and the second equals $\tfrac14$.
\end{itemize}
\end{notesbox}

\begin{practicebox}
Simplify completely and state \emph{all} restrictions. Show the factored line.

\vspace{3pt}
\begin{enumerate}[leftmargin=1.7em, itemsep=8pt]
  \item $\dfrac{x^2+5x+6}{x^2-4}$ \hfill Factored: \ans{$\frac{(x+2)(x+3)}{(x-2)(x+2)}$} \\[4pt]
        Simplest form: \ans{$\frac{x+3}{x-2}$} \quad Restrictions: $x\neq$ \ans{$2,\ -2$}
  \item $\dfrac{9-x^2}{x^2-2x-3}$ \hfill Factored: \ans{$\frac{-(x-3)(x+3)}{(x-3)(x+1)}$} \\[4pt]
        Simplest form: \ans{$-\frac{x+3}{x+1}$} \quad Restrictions: $x\neq$ \ans{$3,\ -1$}
  \item In problem 2, which restriction produces a \emph{hole} and which produces a \emph{vertical
        asymptote}? \ans{$x=3$ cancels $\Rightarrow$ hole; $x=-1$ stays $\Rightarrow$ asymptote}
\end{enumerate}
\end{practicebox}

\begin{teachernote}
\textbf{Pacing.} Sections 1--3 are the lesson; Sections 4--5 are what separates a B from an A. If time
is short, move the monomial example to homework rather than cutting Section~5 --- the equivalence idea
is the A2.EO.1d half of the standard.

\vspace{3pt}
\textbf{The one sentence to repeat.} \emph{Cancelling never restores an input.} Students accept this in
the hook table (they can see $A(3)$ fail) and then forget it two problems later. When they do, send
them back to the $x=3$ column, not to a rule.

\vspace{3pt}
\textbf{Errors to expect.} (1) Cancelling \emph{terms}: $\frac{x^2+9}{x+3}=x+3$, or $\frac{x+3}{3}=x$
--- always disprove numerically, never by assertion. (2) Restrictions read off the \emph{simplified}
denominator, which silently deletes the hole; this is the same mistake that manufactures extraneous
solutions in 4.6, so name the connection now. (3) Cancelling $x-5$ against $5-x$ as if they matched,
losing the $-1$. (4) Writing the simplest form with no restrictions at all --- the answer is the
expression \emph{and} its restriction list. (5) In Section~5, claiming the two expressions are
``different functions'' --- they are equivalent wherever both are defined; they differ only in domain.
\end{teachernote}

\end{document}
